% Metodologia -- Módulo 2 -- Características e sinais auxiliares
\begin{frame}{Módulo 2 -- Características e sinais auxiliares}
  \AipePipelineContext{2}
  \keymessage{As características descrevem a série; a previsão entra como sinal auxiliar, não como critério final de escolha.}
  \vspace{0.25em}

  \begin{columns}[T,onlytextwidth]
    \column{0.53\textwidth}
      \textbf{Características da demanda}
      \begin{itemize}
        \small
        \item ADI e CV$^2$
        \item \textit{Burstiness} e variabilidade
        \item Tendência e sazonalidade
        \item Volume e concentração temporal
      \end{itemize}
    \column{0.47\textwidth}
      \textbf{Sinais auxiliares}
      \begin{itemize}
        \small
        \item Erros de previsão ajudam a medir previsibilidade
        \item O PSE usa sinais para recomendar melhor
        \item A previsão não escolhe a política
      \end{itemize}
  \end{columns}

  \vfill
  \small \textcolor{darkblue}{\textbf{Esses atributos formam $\mathbf{x}_i$, o vetor contextual da série loja-produto.}}
  \note{
    O segundo módulo resume cada série loja-produto em características
    operacionais. ADI e CV ao quadrado capturam intermitência e
    variabilidade, enquanto burstiness, tendência, sazonalidade e volume
    refinam a leitura operacional. Os sinais de previsão entram apenas como
    informação auxiliar sobre previsibilidade. Eles não escolhem a política:
    a escolha continua vindo da avaliação comparativa no simulador e, depois,
    da recomendação aprendida pelo PSE.
  }
\end{frame}
